\subsection{Riesgos operacionales, custodia y adopción}\label{sec:sd8-operacionales}
El registro \path{componentes/escenarios_operacionales_riesgos.csv} enlaza las aplicaciones operacionales con las evaluaciones iniciales, tratamientos y contingencias de T-16. No incorpora dispositivos, funciones ni capacidades ajenas al alcance actual. Rechazo, ausencia de TI, congelamiento y restricciones de maniobra son condiciones conocidas; el evento incierto es que la alternativa elegida no sostenga registro, custodia o adopción. La Tabla~\ref{tab:sd8-siete-caso} permite localizar los siete asuntos expresos del capítulo 19, conservando las reservas de sus IDs.
\begingroup\setlength{\tabcolsep}{3pt}
\begin{longtable}{L{.28\textwidth}L{.16\textwidth}L{.46\textwidth}}
\caption{Riesgos expresos del caso y alternativa sometida a control}\label{tab:sd8-siete-caso}\\\hline
\EncabezadoTabla{Condición conocida} & \EncabezadoTabla{Riesgos T-16} & \EncabezadoTabla{Evento y respuesta}\\\hline\endfirsthead
\hline\EncabezadoTabla{Condición conocida} & \EncabezadoTabla{Riesgos T-16} & \EncabezadoTabla{Evento y respuesta}\\\hline\endhead
Armador rechaza dispositivos. & R-03/10/14 & No informa y el registro pierde salida/regreso. Mantener observación y declaración por personal/radio, sin exigir equipo a bordo; posición sólo con consentimiento expreso y revocable.\\\hline
Apoderado rechaza cámaras. & R-10/11/14 & La alternativa no permite saber quién salió/volvió. Lista, hecho real, bote/monitor y retirada visibles; no cámaras sobre alumnos ni dispositivos en botes menores.\\\hline
Contratista rechaza aplicación. & R-14; R-IN1-01/02 & Acreditación o solicitud queda fuera del registro. Atención asistida con identidad, original y responsable; aplicación voluntaria no condiciona control en portería.\\\hline
Único conocedor del legado. & R-04/12/13 & Ausencia impide interpretar/cerrar o migrar. Transferencia desde mes 6, originales, suplente y ensayo de cierre; sin evidencia bloquear corte y mantener convivencia segura.\\\hline
Corrosión/flexión en pantalán. & R-02/06/11 & Falla montaje, enlace o protección. Aptitud marina, prueba en estructura móvil, repuesto compatible y reposición SD4/T-11; logística/modelos siguen por comprobar.\\\hline
Ventanas canceladas reiteradamente. & R-01; R-IN2-03 & No cabe reprogramación antes de hito. Suspender todos los frentes, consumo único y propagación de red; diez días por campaña no acreditan cobertura global.\\\hline
CLIENTE sin contraparte TI. & R-04/13 & Soporte o validación se asigna a quien no puede ejecutarlos. Synaptix administra/recupera/parchea/vigila; marina valida funcionalmente y designa suplentes, sin presumir disponibilidad.\\\hline
\end{longtable}\endgroup
\FuenteTabla{Caso 07, capítulo 19, restricciones y decisiones 16.1; fichas T-16 y diseños SD3/SD4/SD5/SD7.}

\subsubsection{Custodia marítima, asignación y avisos}
R-03/R-08/R-14 analizan amarra falsamente libre u ocupada por falta de aviso, observación o instrumento. DEC-01 de SD3 y el diseño SD4 distinguen asignación contractual, presencia física y evidencia con antigüedad/calidad; conflicto o dato vencido produce estado indeterminado y verificación humana. No se adjudica el puesto ni se factura permanencia sólo por un sensor o silencio. Se conserva historial y se concilian visitante, contrato, lectura y hecho, sin alterar deuda ni custodia por una inferencia.

R-03/R-10/R-11 consideran regreso cerrado por mensaje tardío, ausencia de cierre o alarma por comunicación retrasada. Hora del hecho y recepción son distintas; una posición no confirma arribo ni el silencio demuestra accidente. La alerta por plan vencido sin cierre se emite en no más de quince minutos y genera revisión/contacto/radio conforme SD3/SD4, sin atribuir búsqueda y salvamento a la plataforma. Operaciones conserva responsabilidad funcional y coordinación con autoridad; se registra la resolución sin inventar un regreso para cerrar alarma. \parencite[decisión 4; numeral 17.4; RT-09.01]{sd8-c07}

R-03/R-08/R-13/R-14 analizan contactos obsoletos, canal caído/saturado y falta de acuse o respuesta a llamada. El caso informa 47\,\% con algún dato obsoleto y 190 de 296 contactados en julio; quedan 106 sin contacto acreditado en ese episodio, no una tasa futura de falla del sistema. Se verifican contactos en renovación y ante rebote/cambio, se conservan versión y origen y se ensayan los tres canales simultáneos: texto, mensajería instantánea y correo. Hasta el último envío a 296 armadores no deben transcurrir más de diez minutos; eso no garantiza entrega, acuse ni ejecución de instrucciones dentro del mismo plazo. La cola y los límites del proveedor deben dimensionarse para ese lote sin sustituirlo por un promedio diario.

El expediente diferencia creación del aviso, envío por canal, respuesta técnica/entrega, acuse humano e instrucción/actuación conocida. Quien no acusa entra en llamada registrada con operador, instante, destinatario y resultado; contacto fallido o nadie disponible para llamar permanece pendiente visible y se escala. La marina puede reforzar amarras por su procedimiento competente, sin afirmar que todos respondieron. Aviso oficial de marejada suspende actividades programadas aunque un armador o proveedor no acuse. \parencite[numerales 7.2 y 9.2; RT-09.01/RT-16.21; capítulo 18, criterio 4]{sd8-c07}

R-03/R-11/R-14 cubren asignación incompatible por eslora, manga, calado, viento o tipo, lista de espera mal aplicada y excepción societaria no autorizada. El 22\,\% en tramo incorrecto y las 90 embarcaciones en espera son diagnóstico, no probabilidad de falla de la nueva regla. Se valida compatibilidad, vigencia de contrato y prioridad versionada; la excepción requiere decisión competente, motivo y auditoría, sin adjudicar automáticamente por calidad de accionista. R-03/R-12 incluyen reserva duplicada, recalada nocturna sin administración y visitante que sale sin hecho facturable: captura local y reserva transaccional conservan ingreso/salida y posterior conciliación contable, sin generar dos cargos. El 9\,\% de recaladas no facturadas identifica exposición histórica. R-03/R-14 revisan documento vencido y legajo erróneo antes de mostrador; el 34\,\% con corrección manual no justifica completar antecedentes inventados ni invadir autorización marítima. Sin prueba de vigencia se mantiene la tarea pendiente y la autoridad decide zarpe.

\subsubsection{Escuela y registro en condiciones reales}
R-03/R-10/R-11/R-14 relacionan retirada anticipada invisible, llegada tardía, cambio de bote/monitor y regreso sin cierre con el episodio de cuarenta minutos y los 470 menores anuales. Autorización previa, asistencia real, salida, cambio, regreso y entrega son hechos distintos. Escuela/portería conservan la misma identidad y versión; una retirada no puede dejar al alumno como en el agua ni una matrícula confirmar presencia. La clase verifica lista efectiva y autorización antes de salir; dato de emergencia ausente o estado incierto bloquea la salida/entrega dependiente y activa contraste humano. Se conserva registro de quién comprobó y resolvió, sin usar cámara sobre alumnos como contingencia.

Monitores mojados o con manos ocupadas pueden no completar captura por guantes, brillo, caída, batería o cobertura. SD3 y SD4 seleccionan el Samsung Galaxy XCover7 Pro como terminal de referencia, conforme a T-11, ítems 61--62. Su variante local y aptitud para cada perfil requieren verificación; la referencia del modelo no acredita el ensayo de uso. SD4/T-11 especifican sujeción, batería y pruebas de turno; dos interacciones/diez segundos para salida/regreso y sesenta segundos para una clase de hasta dieciocho botes se verifican por tarea, sin obligar a tocar pantalla durante maniobra. Las autonomías de doce horas en terminal y ocho en Ilque se prueban en su uso real, no se deducen de una ficha comercial. Ante falla se protege la vigilancia/custodia y se utiliza el procedimiento BC-03 con folio y posterior conciliación; el papel temporal no acredita disponibilidad electrónica.

R-10 aplica visibilidad mínima de salud y datos personales con autorización de finalidad, cifrado y acceso auditado de SD4/SD5. El monitor sólo accede al antecedente que requiere su actividad. R-03/R-10/R-13 analizan aviso al apoderado omitido, tardío o asociado a otro alumno: se valida alumno/destinatario/contacto y versión, se registra evento de origen y estado de envío/acuse; la urgencia usa llamada/contacto humano y no difunde salud en un mensaje general. Al rehabilitar se concilian pendientes sin enviar duplicados ni borrar el aviso equivocado.

\subsubsection{Faenas, combustible, ambiente y boyas}
R-04/R-11/R-14 evalúan pérdida del saber de eslingado, ficha incorrecta y operador nuevo. El caso registra cero de 296 planes documentados y conocimiento en dos personas. Varadero conserva plan validado, versión, embarcación y condiciones antes de izar; una ficha informática no certifica mecánica ni habilita a un operador. El plan se consulta antes de maniobra y la documentación se completa en condición segura con embarcación en calzos; no se exige interacción durante carga suspendida o amarre. Si falta plan válido o competencia, se suspende la maniobra dependiente y se obtiene validación técnica competente.

R-01/R-11/R-14 contemplan cancelación meteorológica, orden físico incompatible en explanada y fallo del único travelift. La agenda conserva posición, obstáculos, predecesoras y disponibilidad del equipo: reprogramar una botadura no permite atravesar otra embarcación ni recuperar una ventana suspendida. Se verifica acceso y condición segura antes de reanudar; reparación mecánica es actividad competente externa, no desarrollo de software. R-14/R-IN1-01 consideran empresas con acreditación vencida o formación ambiental no probada: sólo 41\,\% de 62 empresas estaban completas y unas 140 personas externas requieren control. Portería valida persona/empresa/faena/vigencia y formación antes de acceso; la negativa a instalar aplicación activa registro asistido, sin levantar restricciones ni inventar capacitación aprobada.

R-03/R-11/R-15 cubren mezcla de residuos sin generador, retirada sin gestor/evidencia y derrame sin acta/seguimiento. Siete derrames menores y dos actas revelan déficit probatorio. Se registra generador, recipiente/lote, tipo, traslado, gestor, documento y discrepancia; la entrega incompleta queda pendiente y no se declara disposición conforme. Un derrame activa seguridad/contención competente y acta con seguimiento; la plataforma aporta trazabilidad, sin afirmar que haya ejecutado limpieza o certificación.

R-03/R-08/R-15 y R-IN5-01 cubren medidor mal asociado, lectura ausente/alterada y calibración/evidencia insuficientes. Cero de 320 amarras medidas y 38\,\% de costo no recuperado son antecedentes del caso, no ahorro garantizado. SD5 conserva asociación histórica medidor/ocupante/período, unidad, calidad y rectificación; dato no validado impide publicar cobro o reducción ambiental atribuible. Se comprueba medición y se concilia, sin repartir por una fórmula no aprobada ni imputar cero a una lectura faltante. R-03/R-11 evalúan despacho manual/automático incorrecto, prueba perdida y equipo en zona clasificada: sólo se habilita variante/interfaz acreditada, conserva totalizador y evidencia y se suspende despacho inseguro; un terminal general no pasa a ser apto para surtidor por esta contingencia.

R-02/R-03/R-11/R-13 evalúan ocupación de boya vieja, inspección no registrada y reposición impedida por mar. Las cuarenta boyas a seis millas carecen de energía, datos y personal; la inspección prevista dos veces por semana no prueba ocupación continua. Se conserva fecha/autor/calidad y se muestra incertidumbre hasta verificar, con captura local y sincronización al volver. El traslado requiere embarcación y ventana segura; no se fuerza salida para cumplir un reloj ni se inventa una reparación o visita realizada. SD4 propone repuestos precargados adicionales al instalado, recepción marina y plan de vida útil; su compromiso de reemplazo crítico en cuatro horas desde diagnóstico y logística real por modelo/sitio deben comprobarse, especialmente Ilque. Corrosión/flexión y reposición reiterada mantienen abierta esa dependencia hasta acreditar variante, stock y tiempo de sustitución.

\subsubsection{Adopción y continuidad de compromisos}
R-14 considera vuelta a cuadernos, doble registro y abandono por fricción, apoyo insuficiente o formación tardía. La contingencia asistida conserva identidad/folio y conciliación; no se acepta el cuaderno como sistema oficial permanente. Más de la mitad del personal estacional cambia anualmente y los monitores se renuevan por completo: formación por perfil antes de temporada y altas, evaluación individual y suplencia de SD7 deben preceder la habilitación, fuera del congelamiento. Una plaza ofrecida no es una persona capacitada ni reemplaza práctica de tarea.

R-04/R-05/R-13 evalúan referente funcional ausente, dedicación insuficiente y cola de validaciones que la marina no puede absorber. La dedicación funcional se establece mediante agenda y suplente por proceso; Synaptix conserva administración técnica, mesa, guardia y coordinación. Si falta revisor o capacidad se bloquea decisión/corte dependiente y se prioriza recuperación sin aceptación tácita. R-13 incluye falta de inglés en temporada, noche/emergencia, traslado y reemplazo tardíos y degradación durante 36 meses: T-15 distingue mesa estacional bilingüe, guardia 24\,$\times$\,7 y especialistas, pero los cuadrantes/capacidad y logística no se acreditan con nombres propuestos. La llamada de emergencia sigue procedimiento competente aun cuando no haya mesa ordinaria.

R-15/R-IN5-01/02 relacionan datos ambientales insuficientes y pérdida de certificación/cliente charter antes de temporada 2029. Las tres no conformidades son residuos por generador/destino, consumo real por amarra y formación ambiental de contratistas. Portal, inglés y consumo probado también condicionan al charter, que representa 22\,\% de ingresos por amarra, no de todos los ingresos. Medición temprana, archivo reconstruible y hitos ambientales tienen evidencia propia; un piloto o la instalación de medidores no demuestran cierre de auditoría ni conservación del contrato. Resultado insuficiente se comunica como tal y activa corrección en ventana válida, sin inventar certificación concedida.

R-03/R-10/R-11 analizan disputa por cobro, excepción societaria o acceso bloqueado por mora. Finanzas mantiene conciliación y expediente, y una restricción al acceso a embarcación sólo se aplica con decisión formal fundada y registrada de la compañía; no existe bloqueo automático por deuda ni se atribuye facultad judicial a Synaptix. Los expedientes de abandonadas conservan preservación y decisiones competentes. Las cinco innovaciones vigentes mantienen sus 18 fichas R-IN1--R-IN5: participación/actualización y atención asistida IN1, cohorte/contratista/ventana e interfaz IN2, revisión/permiso/continuidad IN3, sesgo de incentivo/cierre IN4 y medición/adhesión/privacidad IN5. Sus mitigaciones y contingencias de SD13/T-19 se siguen por esos IDs, sin sustituir servicio base por adopción voluntaria ni crear una segunda reserva. \parencites[arts.~29, 85 y 89; pp.~19--20, 44--45]{sd8-ba}[RT-26.04; p.~44]{sd8-bt}[capítulos 1, 19; decisiones 17/18; restricción 15]{sd8-c07}
